\subsection{APPLICATIONS: I. RATIONAL INTEGERS}

Consider the commutative monoid \(\mathbf{N}\) of natural numbers with law of composition addition; all the elements of \(\mathbf{N}\) are cancellable under this law (Set Theory, III, § 5, no. 2, Corollary 3). The group of differences of \(\mathbf{N}\) is denoted by \(\mathbf{Z}\) ; its elements are called the rational integers; its law is called addition of rational integers and also denoted by \(+\) . The canonical homomorphism from \(\mathbf{N}\) to \(\mathbf{Z}\) is injective and we shall identify each element of \(\mathbf{N}\) with its image in \(\mathbf{Z}\) . The elements of \(\mathbf{Z}\) are by definition the equivalence classes determined in \(\mathbf{N} \times \mathbf{N}\) by the relation between \((m_1, n_1)\) and \((m_2, n_2)\) which is written \(m_1 + n_2 = m_2 + n_1\) ; an element \(m\) of \(\mathbf{N}\) is identified with the class consisting of the elements \((m + n, n)\) , where \(n \in N\) ; it admits as negative in Z the class of elements \((n, m + n)\) . Every element \((p, q)\) of \(N \times N\) may be written in the form \((m + n, n)\) if \(p \geqslant q\) or in the form \((n, m + n)\) if \(p \leqslant q\) ; it follows that Z is the union of N and the set of negatives of the elements of N. The identity element 0 is the only element of N whose negative belongs to N.

For every natural number m, -m denotes the negative rational integer of m and -N denotes the set of elements -m for \(m \in N\) . Then

\[
\mathbf {Z} = \mathbf {N} \cup (- \mathbf {N}) \quad \text { and } \quad \mathbf {N} \cap (- \mathbf {N}) = \{0 \}.
\]

for \(m \in \mathbf{N}\) , \(m = -m\) if and only if \(m = 0\) .

Let \(m\) and \(n\) be two natural numbers;

\begin{enumerate}
\def\labelenumi{(\alph{enumi})}
\item
  if \(m \geqslant n\) , then \(m + (-n) = p\) , where \(p\) is the element of \(\mathbf{N}\) such that \(m = n + p\) ;
\item
  if \(m \leqslant n\) , then \(m + (-n) = -p\) , where \(p\) is the element of \(\mathbf{N}\) such that \(m + p = n\) ;
\end{enumerate}

\[
(c) (- m) + (- n) = - (m + n).
\]

Properties (b) and (c) follow from no. 3, Proposition 4; as \(\mathbf{Z} = \mathbf{N} \cup (-\mathbf{N})\) , addition in \(\mathbf{N}\) and properties (a), (b) and (c) describe completely addition in \(\mathbf{Z}\) .

More generally -x is used to denote the negative of an arbitrary rational integer x; the composition \(x + (-y)\) is abbreviated to x - y (cf.~no. 8).

The order relation \(\leqslant\) between natural numbers is characterized by the following property: \(m \leqslant n\) if and only if there exists an integer \(p \in N\) such that \(m + p = n\) (Set Theory, III, §3, no. 6, Proposition 13 and §5, no. 2, Proposition 2). The relation \(y - x \in N\) between rational integers x and y is a total order relation on Z which extends the order relation \(\leqslant\) on N. For, for all \(x \in Z\) , \(x - x = 0 \in N\) ; if \(y - x \in N\) and \(z - y \in N\) , then

\[
z - x = (z - y) + (y - x) \in \mathbf {N},
\]

for \(\mathbf{N}\) is stable under addition; if \(y - x \in \mathbf{N}\) and \(x - y \in \mathbf{N}\) , then \(y - x = 0\) , for \(0\) is the only element of \(\mathbf{N}\) whose negative belongs to \(\mathbf{N}\) ; for arbitrary rational integers \(x\) and \(y, y - x \in \mathbf{N}\) or \(x - y \in \mathbf{N}\) , for \(\mathbf{Z} = \mathbf{N} \cup (-\mathbf{N})\) ; finally, if \(x\) and \(y\) are natural numbers, then \(y - x \in \mathbf{N}\) if and only if there exists \(p \in \mathbf{N}\) such that \(x + p = y\) . This order relation is also denoted by \(\leqslant\) .

Henceforth when Z is considered as an ordered set, it will always be, unless otherwise mentioned, with the ordering that has just been defined, the natural numbers are identified with the integers \(\geqslant 0\) ; they are also called positive integers; the integers \(\leqslant 0\) , negatives of the positive integers, are called negative integers; the integers \textgreater{} 0 (resp. \textless{} 0) are called strictly positive (resp. strictly negative); the set of integers \textgreater{} 0 is sometimes denoted by N*.

Let \(x, y\) and \(z\) be three rational integers; then \(x \leqslant y\) if and only if \(x + z \leqslant y + z\) . For \(x - y = (x + z) - (y + z)\) . This property is expressed by saying that the order relation on Z is invariant under translation.

